\documentclass{article}
\usepackage{graphicx} 
\usepackage{amsfonts}
\usepackage{amsmath}
\usepackage[bookmarksdepth=3]{hyperref}  
\title{A Dirty Model for Multi task Learning.}
\author{Mohamed Amine Chalhy}




\begin{document}

\date{}
\maketitle
\tableofcontents  

\newpage

\section{Problem and introduction}
As introduced in the Dirty Model by Jalali et al.~\cite{jalali2010dirty}

In modern science and engineering problems, we often face high-dimensional data
where the number of features ($p$) (variables) is larger than the number of samples ($n$).
Example: ($n=100$) patients, ($p=10,000$) genes.
In such cases, ordinary regression fails, there are infinitely many solutions.
To get meaningful estimates, we must impose structure on the parameters,
that is, we assume something about the form of the true model, such as:
\begin{itemize}
    \item \textbf{Sparsity:} Only a small subset of the $p$ features are actually relevant.
    \item \textbf{Low-rank structure:} The data or parameter matrix lies near a low-dimensional subspace.
    \item \textbf{Sparse Graphical Models:} variables form a sparse dependency network.
\end{itemize}
These structural assumptions make it possible to recover parameters even when ($p > n$).

\subsection{Multi-Task Learning}
In multi-task learning, we have multiple related regression problems (tasks):
$$y^{(1)} = X^{(1)}\bar{\Theta}^{(1)} + \omega^{(1)}, \quad
y^{(2)} = X^{(2)}\bar{\Theta}^{(2)} + \omega^{(2)}, \quad \dots, \quad
y^{(r)} = X^{(r)}\bar{\Theta}^{(r)} + \omega^{(r)}$$
where:
\begin{itemize}
    \item $y^{(k)}$ is the response for the $k$-th task
    \item $r$ = number of tasks (responses)
    \item $X^{(k)} \in \mathbb{R}^{n \times p}$ is the design matrix
    \item $\bar{\theta}^{(k)} \in \mathbb{R}^p$ = regression coefficients for task ($k$)
    \item $\omega^{(k)} \in \mathbb{R}^n$ the noise vector, we assume $\omega^{(k)} \sim \mathcal{N}(0, \sigma^2)$
\end{itemize}

The key idea of Multi-Task Learning is to learn all tasks jointly, exploiting common structure among their parameter vectors.
This can drastically improve estimation accuracy when $n$ is small and $p$ is large.

\subsection{Why do we use a Dirty Model}
In multi-task learning, two popular approaches are elementwise sparse models such as the Lasso, and block-sparse models based on $\ell_1/\ell_q$ regularization.
The Lasso solves each task independently through:
\begin{equation}
\min_{\theta^{(k)}} \frac{1}{2n}\|y^{(k)} - X^{(k)}\theta^{(k)}\|_2^2 + \lambda \|\theta^{(k)}\|_1, \quad k = 1, \dots, r
\end{equation}
This encourages sparsity within each task, but completely ignores potential relationships among tasks.
As a result, when tasks share common features, the Lasso fails to exploit this structure.
To address this, $\ell_1/\ell_q$ block regularization methods were proposed, which jointly estimate all tasks by enforcing shared sparsity patterns across them:
\begin{equation}
\min_{\Theta \in \mathbb{R}^{p \times r}} \frac{1}{2n}\sum_{k=1}^r \|y^{(k)} - X^{(k)}\Theta^{(k)}\|_2^2 + \lambda \sum_{j=1}^p \|\Theta_{j}\|_q
\end{equation}
where $\Theta_j$ denotes the $j$-th row of $\Theta$.
This formulation assumes that the same set of features are active across all tasks enforcing a shared-support structure. However, in realistic scenarios, only some features are shared across tasks, while others are task-specific.
When the degree of overlap is low or the magnitudes of shared coefficients differ substantially, this block-regularized approach can perform worse than Lasso, as shown in prior studies.

To overcome these limitations, we introduce the Dirty Model, which combines the strengths of both methods.
Instead of assuming a single type of sparsity, the Dirty Model represents the parameter matrix $\Theta$ as a superposition of two components:
$\Theta = B + S$
where:
\begin{itemize}
    \item $B$: captures shared (block-sparse) features across tasks
    \item $S$: captures task-specific (elementwise sparse) features
\end{itemize}
The estimation is performed by solving the following convex optimization problem:
\begin{equation} \label{eq:algorithm}
(\widehat{B}, \widehat{S}) =
\arg\min_{B,S}
\frac{1}{2n}\sum_{k=1}^{r} \|y^{(k)} - X^{(k)}(B^{(k)} + S^{(k)})\|_2^2 + \lambda_b \|B\|_{1,\infty} + \lambda_s \|S\|_{1,1}
\end{equation}
Then we output $\widehat\Theta = \widehat B + \widehat S$.
This formulation allows the model to adapt to any level of feature overlap, leveraging shared information when present while maintaining flexibility for task-specific differences.

\section{Results}
\subsection{Theoretical results}
\subsubsection{Preliminary Information}
Before presenting the main theoretical results, we introduce key definitions and notations that will be used throughout this section.

We stack all the tasks into matrix form:



$$Y = X\bar{\Theta} + W$$
where 
\begin{itemize}
    \item $Y = [y^{(1)}, \dots, y^{(r)}] \in \mathbb{R}^{n \times r}$
    \item $\bar{\Theta} = [\bar{\theta}^{(1)}, \dots, \bar{\theta}^{(r)}] \in \mathbb{R}^{p \times r}$
    \item $W = [\omega^{(1)}, \dots, \omega^{(r)}] \in \mathbb{R}^{n \times r}$
\end{itemize}





Our objective is to estimate the true parameter matrix $\bar{\Theta}$ using the Dirty Model.

Let $\mathcal{U}_k = \operatorname{Supp}(\bar{\theta}^{(k)})$ denote the support of task $k$, and define the overall union of supports as:
$\mathcal{U} = \bigcup_{k=1}^{r} \mathcal{U}_k.$

For each element $\bar{\theta}^{(k)}_j$, we define its sign as:


$$ 
\operatorname{sign}(\bar{\theta}^{(k)}_j) =
\begin{cases}
+1 & \text{if } \bar{\theta}^{(k)}_j > 0, \\
-1 & \text{if } \bar{\theta}^{(k)}_j < 0, \\
0 & \text{if } \bar{\theta}^{(k)}_j = 0.
\end{cases}
$$




An estimator $\hat{\Theta}$ is said to successfully recover the signed support if:
$$\operatorname{sign}(\operatorname{Supp}(\hat{\Theta})) = \operatorname{sign}(\operatorname{Supp}(\bar{\Theta}))$$
This is a stronger condition than recovering the union of active features.
We measure the estimation accuracy using the elementwise $\ell_{\infty}$ norm:
\begin{equation*}
\|\hat{\Theta} - \bar{\Theta}\|_{\infty} =
\max_{j=1,\dots,p} \max_{k=1,\dots,r}
\left| \hat{\theta}^{(k)}_j - \bar{\theta}^{(k)}_j \right|.
\end{equation*}

This metric provides a uniform guarantee over all tasks and features, ensuring that no individual coefficient deviates too much from its true value.

Negahban and Wainwright demonstrated that both the Lasso and the $\ell_1/\ell_\infty$ regularized multi-task models exhibit a phase transition in their ability to recover the true support of the model parameters.

This transition describes a sharp boundary between the regimes of successful and unsuccessful recovery as the number of observations increases.
For the Lasso, the rescaled sample size is defined as:
$$\theta_{\text{Lasso}}(n, p, \alpha) = \frac{n}{s \log(p - s)}$$
where:
\begin{itemize}
    \item $n$ : number of samples
    \item $p$ : total number of features
    \item $s$ : number of non-zero coefficients (true sparsity level)
    \item $\alpha$ : fraction of shared features across tasks
\end{itemize}
Wainwright showed that:
\begin{itemize}
    \item When $\theta_{\text{Lasso}} > 2 + \delta$ for any $\delta > 0$, the Lasso successfully recovers the signed support with probability approaching 1.
    \item When $\theta_{\text{Lasso}} < 2 - \delta$, it fails with probability approaching 1.
\end{itemize}
The term $s \log(p - s)$ represents the effective number of degrees of freedom needed to locate the true sparse signal among all possible combinations of features.

Thus, the Lasso requires the sample size $n$ to grow proportionally to this quantity to guarantee reliable recovery.

In contrast, for the $\ell_1/\ell_{\infty}$ 
regularized multi-task regression, the rescaled sample size is defined as:
$$\theta_{1,\infty}(n, p, \alpha) = \frac{n}{s \log(p - (2 - \alpha)s)}$$
Negahban and Wainwright found that this model exhibits a phase transition at:
$$\theta_{1,\infty} = (4 - 3\alpha)$$

Unlike the Lasso, the $\ell_1/\ell_{\infty}$ formulation enforces shared sparsity across tasks, the same subset of features is assumed active in all tasks.

As a result, the effective dimensionality of the problem is lower when tasks share features $\alpha$, leading to improved performance in that regime.

However, when the overlap between active sets is small (low $\alpha$), this shared support assumption becomes restrictive, and the model can perform worse than Lasso.

The difference in rescaled sample sizes between Lasso and $\ell_1/\ell_\infty$ models therefore arises from their different structural assumptions:
\begin{itemize}
    \item Lasso: assumes independent sparsity for each task.
    \item $\ell_1/\ell_\infty$: assumes joint sparsity across all tasks.
\end{itemize}
Because joint sparsity reduces the number of distinct unknowns, the $\ell_1/\ell_\infty$ model’s effective sample complexity depends not only on $s \log(p - s)$ but also on how many features are shared, expressed through the term $(2 - \alpha)s$.

For the Dirty Model we use the definition of rescaled sample size as $\ell_1/\ell_\infty$ and it achieves a rescaled phase transition at:
$\theta_{\text{Dirty}} = (2 - \alpha)$
which lies strictly before either the Lasso or the $\ell_1/\ell_\infty$ thresholds, except in boundary cases:
\begin{itemize}
    \item $\alpha = 0$ : no sharing $\to$ matches Lasso,
    \item $\alpha = 1$ : full sharing $\to$ matches $\ell_1/\ell_\infty$.
\end{itemize}



\begin{figure}[htbp]
  \centering
  \framebox{\parbox{0.8\textwidth}{\centering
   \includegraphics[width=0.8\textwidth]{thch.png}
  }}
  \caption{Phase Transition Threshold.}
  \label{fig:Phase Transition Threshold}
\end{figure}

\subsubsection{Theorem 1: Deterministic case} (A0)
 $$ \| X^{(k)}_j \|_2 \le \sqrt{2n}, \quad \forall j, k$$
This ensures each column (feature) of the design matrix has a comparable scale.
Without this, features with large norms could dominate and distort regularization.

(A1)
$\gamma_b = 1 - \max_{j \in \mathcal{U}^c} \sum_{k=1}^r
\left|
\left< X^{(k)}_j, X^{(k)}_{\mathcal{U}_k} \right>
\left(\left< X^{(k)}_{\mathcal{U}_k}, X^{(k)}_{\mathcal{U}_k} \right>\right)^{-1}
\right|_1 > 0$

This condition ensures low correlation between irrelevant and relevant features.

(A2)
$C_{\min} := \min_{1 \le k \le r} \lambda_{\min} \left( \frac{1}{n}\left< X^{(k)}_{\mathcal{U}_k}, X^{(k)}_{\mathcal{U}_k} \right>\right) > 0$

This ensures the covariance matrix of the active features is well-conditioned (not singular).

(A3)
$D_{\max} := \max_{1 \le k \le r}
\left\|\left( \frac{1}{n} \left< X^{(k)}_{\mathcal{U}_k}, X^{(k)}_{\mathcal{U}_k}\right> \right)^{-1} \right\|_{\infty, 1}
< \infty$

We don’t want the inverse of the covariance matrix of the active features to explode numerically, which would amplify noise and estimation errors.

Further, we require the regularization penalties be set as
$$\lambda_s > 2(2 - \gamma_s)\sigma \frac{\sqrt{\log(pr)}}{\gamma_s \sqrt{n}}$$
$$\lambda_b > 2(2 - \gamma_b)\sigma \frac{\sqrt{\log(pr)}}{\gamma_b \sqrt{n}}$$

Under A0-A3 and if we set the penalties as above, then with high probability ($\le (1 - c_1 e^{-c_2 n})$), the optimization problem \eqref{eq:algorithm} has a unique optimum $\widehat{\Theta}$ that satisfies:

(a) No False Inclusions + Bounded Error
$$\text{Supp}(\widehat{\Theta}) \subseteq \text{Supp}(\bar{\Theta})$$
and
\begin{equation*}
\|\widehat{\Theta} - \bar{\Theta}\|_{\infty, \infty}
\le
\sqrt{\frac{4\sigma^2 \log(p r)}{n C_{\min}}}\lambda_s D_{\max} =  b_{\min}
\end{equation*}

(b) Exact Signed Support Recovery

If additionally, all true nonzero coefficients are sufficiently large:
$$\min_{(j,k) \in \text{Supp}(\bar{\Theta})} |\bar{\theta}^{(k)}_j| > b_{\min}$$
Then we successfully recover the signed support:
$$\operatorname{sign}(\operatorname{Supp}(\hat{\Theta})) = \operatorname{sign}(\operatorname{Supp}(\bar{\Theta})).$$

Here the positive constants $c_1, c_2$ depend only on $\gamma_s, \gamma_b, \lambda_s, \lambda_b$ and $\sigma$, but are otherwise independent of $n, p, r$, the problem dimensions of interest.

Remark: Condition (a) guarantees that the estimate will have no false inclusions; i.e. all included features will be relevant.

If in addition, we require that it have no false exclusions and that recover the support exactly, we need to impose the assumption in (b) that the non-zero elements are large enough to be detectable above the noise.

\subsubsection{Theorem 2: Random case}
For each task ($k = 1, \dots, r$), the rows $X^{(k)}_i \in \mathbb{R}^p$ are distributed as $X_i^{(k)} \sim \mathcal{N}(0,\Sigma^{(k)})$ where $\Sigma^{(k)} \in \mathbb{R}^{p \times p}$.

We denote by $\Sigma^{(k)}_{\mathcal{V},\mathcal{U}}$ the submatrix of $\Sigma^{(k)}$ with rows indexed by $\mathcal{V}$ and columns indexed by $\mathcal{U}$.

(C1) There is low correlation between irrelevant and relevant features:
$$\gamma_b := 1 - \max_{j \in \mathcal{U}^c} \sum_{k=1}^r
\Big| \Sigma^{(k)}_{j, \mathcal{U}_k} (\Sigma^{(k)}_{\mathcal{U}_k, \mathcal{U}_k})^{-1} \Big|_1 > 0$$

(C2) The covariance matrix of active features is well-conditioned:
$$C_{\min} := \min_{1 \le k \le r} \lambda_{\min}(\Sigma^{(k)}_{\mathcal{U}_k, \mathcal{U}_k}) > 0$$

(C3) The inverse covariance matrices are numerically stable:
$$D_{\max} := \max_{1 \le k \le r} \Big| (\Sigma^{(k)}_{\mathcal{U}_k, \mathcal{U}_k})^{-1} \Big|_{\infty,1} < \infty$$

Let $s := \max_k |\mathcal{U}_k|$.
Then $\lambda_s$ and $\lambda_b$ should be set as:
\begin{equation*}
\lambda_s > \frac{(4\sigma^2 C_{\min} \log(p/r))^{1/2}}{\gamma_s \sqrt{n C_{\min}} - \sqrt{2s \log(p/r)}}
\end{equation*}
and
\begin{equation*}
\lambda_b > \frac{(4\sigma^2 C_{\min} r(r \log(2) + \log(p)))^{1/2}}{\gamma_b \sqrt{n C_{\min}} - \sqrt{2sr(r \log(2) + \log(p))}}.
\end{equation*}

Under C1-C3 and if we set the penalties as above, and if $n$ satisfies:

\begin{equation*}
n > \max\left(\frac{2s \log(p/r)}{C_{\min} \gamma_s^2}, \frac{2sr(r \log(2) + \log(p))}{C_{\min} \gamma_b^2}\right)
\end{equation*}
The optimization problem \eqref{eq:algorithm} satisfies probability at least $1 - c_1 \exp\left(-c_2 (r \log(2) + \log(p))\right) - c_3 \exp(-c_4 \log(r/s))$
for some positive constants ($c_1, \dots, c_4$):

We have same as before
(a) No False Inclusions + Bounded Error
$\text{Supp}(\widehat{\Theta}) \subseteq \text{Supp}(\bar{\Theta})$
and
\begin{equation*}
\|\widehat{\Theta} - \bar{\Theta}\|_{\infty, \infty} \le
\sqrt{\frac{50 \sigma^2 \log(rs)}{n C_{\min}}} +
\lambda_s \left( \frac{4 s}{C_{\min} \sqrt{n}} + D_{\max} \right) = g_{min}.
\end{equation*}

(b) Exact Signed Support Recovery
If additionally, all true nonzero coefficients are sufficiently large:
$$\min_{(j,k) \in \text{Supp}(\bar{\Theta})} |\bar{\theta}^{(k)}_j| > g_{\min}$$

Then we successfully recover the signed support:
$$\operatorname{sign}(\operatorname{Supp}(\hat{\Theta})) = \operatorname{sign}(\operatorname{Supp}(\bar{\Theta})).$$

\subsubsection{Theorem 3: Sharp Phase Transition for 2-Task Gaussian Designs}
This theorem establishes the precise phase transition boundary for the Dirty Model in the special case of two tasks with Gaussian design matrices.
We focus on the case where:
\begin{itemize}
    \item $r = 2$ (two tasks)
    \item Design matrices have rows generated from the standard Gaussian distribution $\mathcal{N}(0, I_{n \times n})$
    \item This ensures conditions C1-C3 are satisfied with $C_{\min} = D_{\max} = 1$
\end{itemize}
We show that the Dirty Model outperforms both Lasso and $\ell_1/\ell_\infty$ block regularization across all $\alpha \in ]0,1[$, matching their performance only at the extremes:
\begin{itemize}
    \item $\alpha = 0$ (no sharing): matches Lasso
    \item $\alpha = 1$ (full sharing): matches $\ell_1/\ell_\infty$
\end{itemize}
For the phase transition analysis, rescaling of the sample size:
For the sharp transition analysis, the regularization parameters must satisfy:

(F1)
\begin{equation*}
\lambda_s > \frac{\left(4\sigma^2(1 - \sqrt{s/n})(\log(r) + \log(p - (2-\alpha)s))\right)^{1/2}}{\sqrt{n} - \sqrt{s} - \sqrt{(2-\alpha)s(\log(r) + \log(p - (2-\alpha)s))}}
\end{equation*}

(F2)
\begin{equation*}
\lambda_b > \frac{\left(4\sigma^2(1 - \sqrt{s/n})r(r\log(2) + \log(p - (2-\alpha)s))\right)^{1/2}}{\sqrt{n} - \sqrt{s} - \sqrt{(1 - \alpha/2)sr(r\log(2) + \log(p - (2-\alpha)s))}}
\end{equation*}
Consider a 2-task regression problem with parameters $(n, p, s, \alpha)$, where the design matrix has rows generated from $\mathcal{N}(0, I_{n \times n})$.

We require that the magnitudes of shared coefficients are approximately balanced across tasks:




$$\max_{j \in \mathcal{B}^*} \left| |\bar{\Theta}^{(1)}_j| - |\bar{\Theta}^{(2)}_j| \right| = o(\lambda_s)$$

where $\mathcal{B}^*$ denotes the set of rows where both entries are non-zero.

Then the estimator $\widehat{\Theta}$ from optimization problem \eqref{eq:algorithm} satisfies:

\textbf{(Success)}
Under F1-F2, and if the rescaled sample size satisfies:
$$\theta(n, p, s, \alpha) > 1$$
then with probability at least $1 - c_1 e^{-c_2 n}$ where $c_1, c_2$ are constants, the estimator $\widehat{\Theta}$ satisfies: (a) and (b) From Theorem 2.

\textbf{(Failure)}
If the rescaled sample size falls below the threshold:
$$\theta(n, p, s, \alpha) < 1$$
then no solution $(\widehat{B}, \widehat{S})$ to the optimization problem \eqref{eq:algorithm} exists for any choices of $\lambda_s$ and $\lambda_b$ such that signed support recovery succeeds.

\subsection{Empirical Results}
In this section, we empirically validate the theoretical predictions of the Dirty Model using controlled experiments on synthetic data.

\subsubsection{Synthetic Data Simulation}
We consider a two-task regression problem as described in Theorem 3, systematically varying the parameters $(n, p, s, \alpha)$ to analyze how the model performs under different structural conditions.

\textbf{Data Generation Protocol:}
We follow the data generation procedure outlined in the original paper, with slight adaptations due to computational constraints:
\begin{itemize}
    \item The design matrices $X$ have entries sampled as $X_{ij} \sim \mathcal{N}(0, 1)$
    \item We fix $p = 128$ and set the sparsity level to $s = \lfloor 0.1p \rfloor$
    \item We explore different overlap ratios $\alpha \in \left\{0.3, \frac{2}{3}, 0.8\right\}$.
    \item For each configuration, the sample size $n$ is varied to control the scaled sample size, as defined in the original study
\end{itemize}
For each $n$, we generate 100 independent problem instances.

In each instance:
\begin{itemize}
    \item The locations of the non-zero entries in the true parameter matrix $\bar{\Theta}$ are chosen uniformly at random, according to $s$ and $\alpha$.
    \item Each non-zero coefficient in $\bar{\Theta}$ is drawn from a uniform distribution over $[-3 \cdot g_{min}, -g_{\min}] \cup [g_{\min}, 3 \cdot g_{min}]$, ensuring that all non-zero coefficients have significant magnitude and well-defined signs.
    \item The response vectors are generated from the linear model with additive Gaussian noise of standard deviation $\Sigma = 0.1$.
\end{itemize}

\textbf{Estimation and Evaluation:}
We estimate $\widehat{\Theta}$ using three models:
\begin{itemize}
    \item The Dirty Model
    \item LASSO
    \item The $\ell_1 / \ell_\infty$ regularized model
\end{itemize}
Due to computational limitations, cross-validation was not performed; instead, the regularization parameters were manually tuned to reproduce the qualitative behaviors reported in the original paper.

\textbf{Success Criterion:}
An instance is considered successful if the estimated coefficient matrix 
$\widehat{\Theta}$ correctly recovers the signed support of the true matrix $\bar{\Theta}$.

To account for small numerical fluctuations, coefficients with absolute values below a tolerance of $10^{-1}$ are ignored.

This success criterion was not explicitly defined in the original paper; it represents a contribution introduced in this work to provide a more practical and reproducible evaluation of signed support recovery under limited computational resources.


We obtained the following empirical results for the different overlap ratios $\alpha$:



\textbf{$\alpha = 0.3$, p = 128}
\begin{figure}[htbp]
  \centering
  \framebox{\parbox{0.8\textwidth}{\centering
   \includegraphics[scale=0.35]{alpha033.png}
  }}
  \caption{Empirical results for $\alpha = 0.3$.}
  \label{fig:alpha023}
\end{figure}

\textbf{$\alpha = \frac{2}{3}$, p = 128}
\begin{figure}[htbp]
  \centering
  \framebox{\parbox{0.8\textwidth}{\centering
   \includegraphics[scale=0.35]{alpha023.png}
  }}
  \caption{Empirical results for $\alpha = 2/3$.}
  \label{fig:alpha023}
\end{figure}

\textbf{$\alpha = 0.8$, p = 128}
\begin{figure}[htbp]
  \centering
  \framebox{\parbox{0.8\textwidth}{\centering
   \includegraphics[scale=0.35]{alpha088.png}
  }}
  \caption{Empirical results for $\alpha =0.8$.}
  \label{fig:alpha023}
\end{figure}

While the results did not exactly reproduce those presented in the original paper—likely due to the absence of cross-validation and computational limitations—the overall trends were consistent with the theoretical predictions:
\begin{itemize}
    \item The Dirty Model consistently achieved the best performance across all settings.
    \item The LASSO model performed relatively well when the overlap ratio $\alpha$ was small, reflecting its strength when tasks share fewer relevant features.
    \item The $\ell_1 / \ell_\infty$ model performed better when $\alpha$ was large, confirming its effectiveness in capturing highly overlapping supports across tasks.
\end{itemize}
These findings reinforce the main theoretical insight that the Dirty Model effectively adapts to varying degrees of task relatedness, outperforming single-structure models across diverse regimes.

\subsubsection{Handwritten Digits Dataset}
We used the Multiple Features dataset from the UCI Machine Learning Repository.
Each feature table was normalized by dividing all values by its maximum value to ensure consistency across features.

The models were trained using different numbers of training samples per class

— n = 10, 20, and 40 — and tested on the remaining ($200-n$) samples per class.


Due to computational limitations, we could not perform full cross-validation. 

Instead, we manually selected parameter values that yielded good performance for both the Dirty Model and the L1/L∞ model.
The chosen parameters were:
\begin{itemize}
    \item $\lambda_s = 0.02$
    \item $\lambda_b = 0.03$
    \item $\lambda_{\text{Lasso}} = 0.001$
    \item $\lambda_{1/\infty} = 0.01$
\end{itemize}
Overall, the Dirty Model performed slightly better when the number of training samples per class was small (n = 10). However, as the number of training samples increased (n = 20 and n = 40), Dirty became slightly less accurate compared to Lasso.
We also observed that the Dirty Model was less stable and took significantly longer to fit, whereas Lasso achieved good performance with minimal parameter tuning.

\bigskip
\hrule
\bigskip

\textbf{Training samples per class: 10}

\texttt{=== Global Errors ===} \\
\texttt{Inf-norm:  0.1026 (10.26\%)} \\
\texttt{Dirty:     0.0937 (9.37\%)} \\
\texttt{Lasso:     0.1184 (11.84\%)}

\texttt{=== Variance of Per-Class Errors ===} \\
\texttt{Inf-norm:  0.002893} \\
\texttt{Dirty:     0.001827} \\
\texttt{Lasso:     0.002871}

\bigskip
\hrule
\bigskip

\textbf{Training samples per class: 20}

\texttt{=== Global Errors ===} \\
\texttt{Inf-norm:  0.0511 (5.11\%)} \\
\texttt{Dirty:     0.0550 (5.50\%)} \\
\texttt{Lasso:     0.0544 (5.44\%)}

\texttt{=== Variance of Per-Class Errors ===} \\
\texttt{Inf-norm:  0.000900} \\
\texttt{Dirty:     0.001256} \\
\texttt{Lasso:     0.000486}


\bigskip
\hrule
\bigskip

\textbf{Training samples per class: 40}

\texttt{=== Global Errors ===} \\
\texttt{Inf-norm:  0.0344 (3.44\%)} \\
\texttt{Dirty:     0.0375 (3.75\%)} \\
\texttt{Lasso:     0.0256 (2.56\%)}

\texttt{=== Variance of Per-Class Errors ===} \\
\texttt{Inf-norm:  0.000424} \\
\texttt{Dirty:     0.000492} \\
\texttt{Lasso:     0.000269}



\bigskip
The results indicate that the Dirty Model achieves slightly better performance when training data is scarce, suggesting that it can better capture useful shared structure across tasks under low-data regimes.

However, as the training set grows (n = 20, 40), Lasso begins to outperform both Dirty and L1/L∞ in terms of global accuracy and per-class stability.
Additionally, the Dirty Model required more computation time and was less stable, whereas Lasso offered competitive results with faster convergence and minimal parameter tuning.

\section{Code Availability}
All Python scripts used for data preprocessing, model training, figure generation, as well as the LaTeX source files for this report, are available at:

\url{https://github.com/Mohemed-Amine-Chalhy/dirty-model-multitask-learning}


\begin{thebibliography}{9}

\bibitem{jalali2010dirty}
Ali Jalali, Sujay Sanghavi, Chao Ruan, and Pradeep Ravikumar.
\newblock A dirty model for multi-task learning.
\newblock \emph{Advances in Neural Information Processing Systems (NeurIPS)}, 2010.

\end{thebibliography}
\end{document}
